\documentclass[12pt,a4paper]{article}
\usepackage[margin=25mm, headheight=15pt, headsep=10pt]{geometry}
\usepackage{fontspec}
\usepackage[table]{xcolor} % Note: table option is required for rowcolors
\usepackage{titlesec}
\usepackage{tocloft}
\usepackage{fancyhdr}
\usepackage{booktabs}
\usepackage{tcolorbox}
\tcbuselibrary{skins, breakable}
\usepackage[colorlinks=true, linkcolor=emerald, urlcolor=emerald, bookmarks=true]{hyperref}

\definecolor{emerald}{RGB}{0,102,51}
\definecolor{darkred}{RGB}{139,0,0}
\definecolor{lightgray}{RGB}{245,245,245}

\usepackage[nil]{babel}
\babelprovide[import, main]{bengali}
\babelprovide[import]{arabic}
\babelfont{rm}[Renderer=HarfBuzz,Scale=1.0]{Noto Serif Bengali}
\babelfont{sf}[Renderer=HarfBuzz]{Noto Sans Bengali}
\babelfont{tt}[Renderer=HarfBuzz]{Noto Sans Bengali}
\babelfont[arabic]{rm}[Renderer=HarfBuzz,Scale=1.1]{Noto Naskh Arabic}

% Section styling
\titleformat{\section}{\Large\bfseries\color{emerald}}{\thesection.}{0.5em}{}
\titleformat{\subsection}{\large\bfseries\color{emerald!85!black}}{\thesubsection.}{0.5em}{}
\titleformat{\subsubsection}{\normalsize\bfseries\color{emerald!70!black}}{\thesubsubsection.}{0.5em}{}
\titlespacing*{\section}{0pt}{18pt}{8pt}

% Fancy Header/Footer setup
\pagestyle{fancy}
\fancyhf{} % clear all header and footer fields
\fancyhead[L]{\color{emerald}\small\textbf{\nouppercase{\leftmark}}}
\fancyfoot[L]{\scriptsize\color{black!50}{\lat{AI}}-সংকলিত, আলেম-যাচাইকৃত নয়}
\fancyfoot[C]{\color{emerald!70!black}\thepage}
\renewcommand{\headrulewidth}{0.4pt}
\renewcommand{\headrule}{\hbox to\headwidth{\color{emerald}\leaders\hrule height \headrulewidth\hfill}}
\renewcommand{\footrulewidth}{0pt}

% Custom boxes
% 1. Ayat/Hadith Box
\newtcolorbox{ayatbox}[1][]{
    enhanced, breakable, 
    colback=emerald!5, 
    colframe=emerald, 
    boxrule=0.5pt, 
    arc=3pt, 
    left=10pt, right=10pt, top=10pt, bottom=10pt,
    #1
}

% 2. Quote/Translation Box
\newtcolorbox{quotebox}[1][]{
    enhanced, breakable,
    frame hidden,
    colback=lightgray,
    borderline west={3pt}{0pt}{emerald},
    left=15pt, right=10pt, top=10pt, bottom=10pt,
    sharp corners,
    #1
}

% 3. AI-generated notice box
\newtcolorbox{warnbox}[1][]{
    enhanced, breakable,
    colback=red!4,
    colframe=darkred,
    boxrule=0.8pt,
    arc=3pt,
    left=10pt, right=10pt, top=10pt, bottom=10pt,
    #1
}

% Commands
\newcommand{\arab}[1]{\foreignlanguage{arabic}{#1}}
\newcommand{\kib}[1]{\textcolor{emerald}{\textbf{#1}}}

% Latin (English) fallback: Noto Serif Bengali has NO Latin glyphs
\newfontfamily\latfont[Renderer=HarfBuzz]{Noto Serif}
\newcommand{\lat}[1]{{\latfont #1}}

\setlength{\parskip}{5pt}
\setlength{\parindent}{0pt}

% Fix for missing bullet in Noto Serif Bengali
\renewcommand{\labelitemi}{$\bullet$}



\begin{document}
\begin{center}{\Large\bfseries\color{emerald} ঘটনা ২: উমার (রাঃ) — জান্নাতের প্রাসাদ পেয়েও নবীজির প্রতি বিনয়}\end{center}
\vspace{1em}
\section*{ঘটনা ২: উমার (রাঃ) — জান্নাতের প্রাসাদ পেয়েও নবীজির প্রতি বিনয়}

\begin{warnbox}
 \textbf{গুরুত্বপূর্ণ নোটিশ (\lat{Important} \lat{Notice}):} এই কন্টেন্টটি \lat{AI} (কৃত্রিম বুদ্ধিমত্তা) দিয়ে প্রস্তুত ও সংকলিত। \lat{AI} কঠোর সূত্র-মান নীতি মেনে শুধু নির্ভরযোগ্য উৎস থেকে তথ্য সংগ্রহ করেছে — কেবল সহীহ/হাসান হাদিস ও সহীহ সনদের আসার রাখা হয়েছে, দুর্বল সনদ বাদ দেওয়া হয়েছে। তবে এটি কোনো আলেম বা মুহাদ্দিস কর্তৃক যাচাইকৃত নয়।
\end{warnbox}


\medskip\hrule\medskip

\subsection*{১. সূত্র কার্ড}

\begin{table}[h]\centering\small
\begin{tabular}{ll}
বিষয় & বিবরণ \\
\hline
\textbf{বর্ণনাকারী} & আবু হুরাইরা (রাঃ) থেকে; সাঈদ ইবনে মুসাইয়্যাব (রহ.) কর্তৃক বর্ণিত \\
\textbf{গ্রন্থ} & সহীহ বুখারি, হাদিস ৩৬৮০; সহীহ মুসলিম, হাদিস ২৩৯৫এ \\
\textbf{অধ্যায়} & কিতাবু ফাযায়িলিস সাহাবা, বাব: উমার ইবনুল খাত্তাবের মর্যাদা \\
\textbf{সনদের মান} & \textbf{সহীহ} (বুখারি ও মুসলিম উভয়ের শর্তানুযায়ী) \\
\textbf{মূল বিষয়} & গীরা (আত্মমর্যাদা) ও নবীজির প্রতি সাহাবির ভালোবাসার গভীরতা \\
\hline
\end{tabular}
\end{table}

\medskip\hrule\medskip

\subsection*{২. পটভূমি (কে, কখন, কোথায়, কেন)}

নবীজি (সাল্লাল্লাহু আলাইহি ওয়া সাল্লাম) সাহাবিদের জান্নাতের \lat{good} \lat{news} (সুসংবাদ) দিতেন এবং তাদের \lat{status} (মর্যাদা) তুলে ধরতেন। \textbf{উমার ইবনুল খাত্তাব (রাঃ)} ছিলেন নবীজির সবচেয়ে প্রিয় সাহাবিদের একজন — ইসলামের দ্বিতীয় খলিফা, সত্য ও \lat{justice} (ন্যায়)-এর প্রতীক।

একদিন নবীজি সাহাবাদের সাথে বসে একটি \textbf{\lat{dream} (স্বপ্ন)-এর} কথা বলছিলেন। নবীদের স্বপ্ন সাধারণ মানুষের \lat{dream} নয় — এটি \lat{revelation} (ওহী)-এর অংশ। নবীজি \lat{dream}-এ জান্নাতের এমন কিছু দেখেছিলেন যা শুনে উমার (রাঃ)-এর চোখে পানি এসে গিয়েছিল।

\medskip\hrule\medskip

\subsection*{৩. পূর্ণ ঘটনার বর্ণনা}

আবু হুরাইরা (রাঃ) বলেন:

\begin{quote}
\textbf{\lat{“}আমরা নবীজির (সাল্লাল্লাহু আলাইহি ওয়া সাল্লাম) কাছে বসা ছিলাম। তিনি বললেন: "ঘুমের মধ্যে আমি নিজেকে জান্নাতে দেখলাম। হঠাৎ একটি প্রাসাদের পাশে দেখলাম, এক মহিলা অযু করছে। আমি জিজ্ঞেস করলাম: 'এই প্রাসাদটি কার?'} \\
\textbf{তারা (ফেরেশতারা) বলল: 'এটি উমার ইবনুল খাত্তাবের।'} \\
\textbf{তখন আমি উমারের গীরা (আত্মমর্যাদা) স্মরণ করলাম এবং (মহিলাটির কাছ থেকে সরে) দ্রুত পেছন ফিরে চলে গেলাম।"} \\
\textbf{উমার (রাঃ) এই কথা শুনে কেঁদে ফেললেন এবং বললেন: "হে আল্লাহর রাসূল! আমার পিতা-মাতা আপনার জন্য উৎসর্গ হোক! আপনার প্রতি আমি কি গীরা (ঈর্ষা/আত্মমর্যাদা) পোষণ করব? আপনার প্রতি তো আমার ঈর্ষা থাকতেই পারে না! আপনি তো আমার মনিব ও পথপ্রদর্শক!"\lat{”}}
\end{quote}

\medskip\hrule\medskip

\subsection*{৪. শব্দের ব্যাখ্যা}

\begin{table}[h]\centering\small
\begin{tabular}{ll}
শব্দ & অর্থ \\
\hline
\textbf{গীরা (\arab{غيرَة})} & আত্মমর্যাদা; নিজের অধিকার, পরিবার বা সম্মান নিয়ে যত্নবান থাকা। ইসলামে স্বামীর গীরা প্রশংসনীয়। \\
\textbf{তাহারাত / অযু} & পবিত্রতা অর্জন — জান্নাতে এটি শারীরিক ময়লা দূর করার জন্য নয়, বরং সৌন্দর্য বৃদ্ধির জন্য \\
\textbf{ফেদা} & \lat{“}আমার পিতা-মাতা আপনার জন্য উৎসর্গ\lat{”} — আরবি সংস্কৃতিতে চরম ভালোবাসা ও সম্মান প্রকাশের বাক্য \\
\hline
\end{tabular}
\end{table}

\textbf{গীরা ও কিবরের পার্থক্য:}
\begin{itemize}
\item \textbf{কিবর (অহংকার):} নিজেকে অন্যের চেয়ে বড় মনে করা, মানুষকে হেয় করা — সম্পূর্ণ নিন্দনীয়।
\item \textbf{গীরা (আত্মমর্যাদা):} নিজের পরিবার-সম্মান রক্ষায় সজাগ থাকা — প্রশংসনীয়, নবীজি একে ভালোবাসতেন।
\item উমারের গীরা ছিল নিন্দনীয় ঈর্ষা নয়, বরং সতীত্ব ও সম্মানের স্বাভাবিক মমতা — তাই নবীজি তাকে প্রশংসা করেছেন।
\end{itemize}
\medskip\hrule\medskip

\subsection*{৫. সংশ্লিষ্ট হাদিস ও আয়াত}

\textbf{হাদিস (বুখারি ৫২২৬):} নবীজি (সা.) বলেছেন: \textbf{\lat{“}আল্লাহ তো গীরা পছন্দ করেন... আর মুমিনও গীরা করে।\lat{”}} — প্রমাণ করে, গীরা নিন্দনীয় নয়।

\textbf{হাদিস (বুখারি ৩৬৮২):} নবীজি বলেছেন: \textbf{\lat{“}আমার উম্মতের মধ্যে উমারের মতো কেউ নেই যে শয়তান তাকে পথ থেকে ফিরিয়ে দেয় — যখনই সে (উমার) কোনো পথে হাঁটে, শয়তান অন্য পথে চলে যায়।\lat{”}}

\textbf{আয়াত (সূরা আত-তাওবা ৯:৪০):}
\begin{quote}
\textbf{\lat{“}তোমরা তাকে সাহায্য না করলে (জেনে রেখো), আল্লাহ তো তাকে সাহায্য করেছেন... তিনি তার সাথীকে (আবু বকরকে) বলেছিলেন: দুঃখ করো না, আল্লাহ আমাদের সঙ্গে আছেন।\lat{”}}
\end{quote}

\textbf{আয়াত (সূরা আল-হাশর ৫৯:৯):} আনসারদের প্রশংসায় — \textbf{\lat{“}যারা তাদের আগে (মুহাজিরদের) ঘরে স্থান দিয়েছে এবং ঈমান এনেছে, তারা মুহাজিরদের ভালোবাসে।\lat{”}}

\medskip\hrule\medskip

\subsection*{৬. রেফারেন্স}

\begin{itemize}
\item সহীহ বুখারি, হাদিস ৩৬৮০ (কিতাবু ফাযায়িলিস সাহাবা)
\item সহীহ মুসলিম, হাদিস ২৩৯৫এ (ফাযায়িলুস সাহাবা)
\item সহীহ বুখারি, হাদিস ৫২২৬ ও ৩৬৮২
\item ফাতহুল বারি (ইবনে হাজার), ৭/৫৭ — উমারের কান্নার কারণ নিয়ে ব্যাখ্যা
\end{itemize}

\end{document}
